% Chapter1/chapter1.tex -- Introduction
% -----------------------------------------------------------------------------
% The introduction serves three purposes:
%   1. Motivate the problem (why does it matter?)
%   2. State your objectives (what specifically did YOU do?)
%   3. Orientate the reader (what does the rest of the report contain?)
%
% Keep Section 1.1 broad and accessible -- your examiner may not be an ML expert.
% Keep Section 1.2 specific and measurable -- these become your evaluation criteria.

\chapter{Introduction}
\label{ch:introduction}

\section{Motivations}
\label{sec:motivations}

Large language models (LLMs) have demonstrated remarkable capabilities
across a wide range of natural language understanding and generation tasks,
but their deployment in resource-constrained environments remains a
significant challenge \cite{PLACEHOLDER}. State-of-the-art models such as
GPT-4 and LLaMA require hundreds of gigabytes of memory and specialised
GPU clusters to achieve practical inference throughput, making them
inaccessible for embedded, edge, and real-time applications.

TinyLlama-1.1B \cite{PLACEHOLDER} represents a class of compact LLMs that
seek to balance capability with deployability. With only 1.1 billion
parameters -- roughly two orders of magnitude smaller than frontier models --
TinyLlama retains competitive performance on common benchmarks while fitting
within a much smaller memory footprint. Nevertheless, even at this scale,
running inference on a conventional CPU is far too slow for real-time use,
and GPUs consume too much power for battery-operated or always-on deployments.

Field-programmable gate arrays (FPGAs) offer a compelling alternative.
Unlike fixed-function ASICs, FPGAs are reconfigurable: the hardware can be
tailored precisely to the computational patterns of a specific model, enabling
high-throughput, low-latency inference at a fraction of the power cost of a GPU
\cite{PLACEHOLDER}. Recent advances in high-level synthesis (HLS) tools --
particularly Xilinx Vitis HLS -- have substantially lowered the barrier to
FPGA development by allowing engineers to describe hardware behaviour in C++
and have the tool generate the corresponding RTL automatically.

\section{Objectives and Scope}
\label{sec:objectives}

This project aims to implement and optimise the inference pipeline of
TinyLlama-1.1B on the Xilinx ZCU102 evaluation board using Vitis HLS and
the Xilinx Runtime (XRT) framework. The primary objectives are as follows.

\begin{enumerate}
  \item Implement the core transformer components of TinyLlama -- including
        RMSNorm, matrix--vector multiplication, multi-head attention (MHA),
        and the SwiGLU feed-forward network -- as optimised HLS kernels.

  \item Apply INT8 post-training quantisation (PTQ) to reduce memory bandwidth
        requirements and exploit the ZCU102's DSP blocks efficiently, while
        maintaining acceptable output quality.

  \item Apply HLS optimisation directives (PIPELINE, UNROLL, ARRAY\_PARTITION)
        to minimise the initiation interval and achieve the target throughput
        within the ZCU102's resource budget.

  \item Evaluate the final implementation against a CPU software baseline in
        terms of inference throughput (tokens per second), latency per token,
        and energy efficiency (tokens per joule).
\end{enumerate}

The scope of this project is limited to autoregressive token generation
(i.e., the decode phase of inference). Prefill-phase optimisation and
training are out of scope.

\section{Organisation of Report}
\label{sec:organisation}

The remainder of this report is organised as follows.
\Cref{ch:litreview} reviews the relevant literature on FPGA-based
transformer acceleration, INT8 quantisation techniques, and related HLS
optimisation strategies. \Cref{ch:architecture} presents the overall system
architecture and the design decisions made at the hardware level.
\Cref{ch:implementation} describes the detailed implementation of each HLS
kernel and the software driver layer. \Cref{ch:results} presents the
experimental results including resource utilisation, performance benchmarks,
and comparison against baseline systems. \Cref{ch:conclusion} draws
conclusions and identifies directions for future work.
